\hypertarget{chapter-49-distributed-c}{%
\section{CHAPTER 49: DISTRIBUTED C}\label{chapter-49-distributed-c}}

\hypertarget{distributed-c}{%
\section{DISTRIBUTED C++}\label{distributed-c}}

Moving beyond a single process: Networking, RPC, and Consensus.

\hypertarget{serialization-the-box-and-label-problem}{%
\subsubsection{1. Serialization: The ``Box and Label''
Problem}\label{serialization-the-box-and-label-problem}}

When you send an object (like a \passthrough{\lstinline!User!} class)
over the network, you can't just send the memory address. The address
\passthrough{\lstinline!0x123!} on your computer doesn't mean anything
to another computer across the world.

Instead, you have to \textbf{Serialize} it. This is like taking a LEGO
castle, breaking it down into individual bricks, putting them in a
numbered box with instructions, and shipping it. The receiver then
\textbf{Deserializes} itrebuilding the castle brick-by-brick.

\hypertarget{serialization-binary-protocols}{%
\paragraph{1.1 Serialization (Binary
Protocols)}\label{serialization-binary-protocols}}

Efficiently packing data for network transmission.

\begin{lstlisting}[language={C++}]
#include <vector>
#include <cstring>
#include <string>

// Simple Binary Serializer
class Buffer {
    std::vector<uint8_t> data;
public:
    template<typename T>
    void write(const T& val) {
        static_assert(std::is_trivially_copyable_v<T>);
        const uint8_t* ptr = reinterpret_cast<const uint8_t*>(&val);
        data.insert(data.end(), ptr, ptr + sizeof(T));
    }

    void write_string(const std::string& s) {
        write<uint32_t>(s.size());
        const uint8_t* ptr = reinterpret_cast<const uint8_t*>(s.data());
        data.insert(data.end(), ptr, ptr + s.size());
    }

    const uint8_t* begin() const { return data.data(); }
    size_t size() const { return data.size(); }
};
\end{lstlisting}

\hypertarget{rpc-remote-procedure-call-concept}{%
\subsubsection{16.2 RPC (Remote Procedure Call)
Concept}\label{rpc-remote-procedure-call-concept}}

Calling a function on another machine.

\textbf{Stub Interface:}

\begin{lstlisting}[language={C++}]
// User Code
// auto result = service.Add(5, 3);

// Generated Stub
int Add(int a, int b) {
    Buffer buf;
    buf.write(101); // Function ID for 'Add'
    buf.write(a);
    buf.write(b);
    return network.send_and_wait(buf); // Blocks
}
\end{lstlisting}

\hypertarget{consensus-raft-basics}{%
\subsubsection{16.3 Consensus (Raft
Basics)}\label{consensus-raft-basics}}

Distributed systems need to agree on state.

\textbf{Raft State Machine:}

\begin{lstlisting}[language={C++}]
enum class State { Follower, Candidate, Leader };

struct Node {
    State state = State::Follower;
    int current_term = 0;
    int voted_for = -1;

    void on_timeout() {
        if (state == State::Follower) {
            state = State::Candidate;
            current_term++;
            voted_for = my_id;
            request_votes();
        }
    }
};
\end{lstlisting}
